\clearpage\setcounter{chapter}{2}
\chapter{Hilbert 空间上无界算子的基本概念}\label{chap:unbounded-operators}

偏微分方程及其数值离散化的许多现代理论, 在 Hilbert 空间及其上算子的框架中研究最为简洁优美. 我们假定读者熟悉 Hilbert 空间及其上有界线性算子的基本定义和结论. 这里着重讨论闭的无界线性算子, 特别关注伴随与对偶. 这同样是一套成熟且发展完善的理论, 尽管受到的研究较少. 本章大多数结论可以在许多参考文献中找到, 如 [\BookNumber{23}, Chapter 2.6], [\BookNumber{56}, Chapter III, Section 5 及 Chapter IV, Section 5.2], 或 [\BookNumber{71}, Chapter II, Section 6 及 Chapter VII]; 这些文献常在更一般的 Banach 空间甚至更一般的拓扑向量空间中讨论. 不过, 对我们的应用而言, Hilbert 空间的背景已经足够, 而且还能使论述简短而较为自足. 在本章末尾, 我们将回到上一章中作为 de Rham 复形微分出现的微分算子, 将它们解释为 Hilbert 空间上的闭无界算子, 并计算其伴随. 主要内容概括于本章末尾的 Box~\ref{box:unbounded-operators-summary}.
